\section{Limitations}

While our results establish the viability of super-nominal payload manipulation on fixed/given hardware, several aspects warrant further discussion. \textbf{Rigid Attachment Choice.} Even though the rigid attachment experimental choice simplifies analysis by eliminating offset moments, it introduces two critical limitations in real-world execution (see accompanying video): dynamic effects from payload oscillation when objects are not perfectly rigid, and unmodeled torques when the object's center of mass is significantly offset from the end-effector's origin. These assumptions can lead to trajectory tracking errors or even task failures, particularly when handling flexible materials or asymmetric objects. Looking ahead, extending our approach to objects with internal dynamics will require explicit robot-object dynamics models to capture effects such as rotational inertia, slippage, and uneven mass distributions. Incorporating these nonlinearities, together with friction modeling for fingered grippers and suction force modeling for suction-based end-effectors, remains an important direction. Moreover, expanding the conditioning space to whole-body contact forces, fingertip wrench cones, and suction dynamics, as well as embedding more advanced controllers (e.g., impedance or admittance), will allow the same ``constraint-aware by design'' principle to scale to interaction-rich manipulation tasks. \textbf{Test-Time Robustness.} Assessing test-time robustness requires long-term deployment studies and, in many cases, proprietary Original Equipment Manufacturer (OEM) data on wear-and-tear. Whether robots should operate continuously at or near maximum capacity is also a design question for robot OEMs. While our method generates trajectories that respect manufacturer-specified torque limits, joints do run hotter under high payloads, motivating further study of safety margins under sustained operation. \textbf{Embodiment Generalization.} We acknowledge that embodiment understanding must be more directly incorporated into the training process to enable generalization across different robotic platforms. To this end, we are developing more general embodiment-constrained conditioning mechanisms as part of ongoing work.

Our work also raises important theoretical questions about completeness and optimality guarantees for learning-based motion planning. Even though we can argue that diffusion models trained on optimal trajectories will generate similarly optimal solutions, establishing formal completeness guarantees remains an open challenge.
Moreover, while our method demonstrates advantages in inference speed, traditional optimization-based approaches offer clearer constraint satisfaction guarantees and greater flexibility in adapting to new constraints without retraining. That said, training is not a bottleneck: we generate training trajectories across multiple payload conditions, with each trajectory verified for kinematic and dynamic feasibility throughout the robot's operational workspace. The compact dimensionality of our denoised samples enables fast retraining of the UNet and efficient creation of a high-quality dataset with comprehensive operational space coverage, but without the need for real-world payload manipulation data (assuming the existence of a robot dynamics model, which also lends to system modularity).

More rigorous evaluations comparing CPU versus GPU performance and Python versus C++ implementations would enable direct baseline comparisons but standardizing these benchmarks remains challenging. Our evaluation of baseline methods maintains industry-standard implementations of baselines without additional optimizations for deployment infrastructure.
Several other under-explored areas of research remain: optimal training data selection for maximizing planner generalization, disambiguating data imbalance from model generalization to explain the high variance in torques for higher payloads, leveraging diffusion models' continual learning capabilities at industrial scale, and conditioning on real-world sensor modalities like point clouds for deployment.